\section{Data and Methods}\label{sec:data-methods}

\subsection{Experimental Solubility and Speciation Data}
The \COtwo-solubility data set contains 161 equilibrium partial-pressure records compiled from six VLE sources\cite{Aronu2011,Hilliard2008,Idris2014,Jou1995,Mamun2005,Xu2011}.  Each canonical record retains a stable source key, source filename, and one-based source-row identity in addition to amine mass fraction, temperature, \COtwo loading, and measured \COtwo pressure.  The aggregate is generated from a row-level inclusion manifest; six Xu records carry explicit nominal-temperature normalization from 100.5--\SI{101.1}{\degreeCelsius} to \SI{100}{\degreeCelsius} or 120.4--\SI{121.8}{\degreeCelsius} to \SI{120}{\degreeCelsius}.  The speciation data set contains 74 NMR-derived liquid-composition records compiled from three spectroscopic studies\cite{Bottinger2008,Jakobsen2005,Matin2012}.  The speciation records are expressed on a common loading basis and reconciled to satisfy carbon, amine, water, and charge balances before model comparison.

\begin{center}
    \begin{minipage}{\linewidth}
    \centering
    \captionsetup{hypcap=false}
    \captionof{table}{Experimental evidence base used for the \MEA \COtwo-solubility and speciation evaluation.}
    \label{tab:mea-data-inventory}
    \small
    \begin{tabularx}{\linewidth}{lclXX}
        \toprule
        Data family & Records & Primary sources & Conditions & Observables and residual role \\
        \midrule
        VLE & 161 & \cite{Aronu2011,Hilliard2008,Idris2014,Jou1995,Mamun2005,Xu2011} &
        30 wt\% \MEA, 40--\SI{120}{\degreeCelsius}, \(\alpha=0.003\)--0.689, \(P_{\COtwo}=1.5\times10^{-3}\)--\(2.804\times10^3\) kPa &
        Measured \COtwo pressure at fixed loading; logarithmic solubility residual \\
        Speciation & 74 & \cite{Bottinger2008,Jakobsen2005,Matin2012} &
        30 wt\% \MEA, 20--\SI{80}{\degreeCelsius}, \(\alpha=0.064\)--0.990 &
        True-species mole fractions; direct-positive, reported-zero, and balance-inferred target roles \\
        \bottomrule
    \end{tabularx}
    \end{minipage}
\end{center}

\subsection{Species Basis and Parameter Hierarchy}
The activity-based calculation uses the nine-species liquid basis defined in \cref{eq:true-species-vector}.  The parameter hierarchy separates inherited molecular parameters, SSM+DS ion conventions from modified-Born ePC-SAFT literature, provisional fixed inputs for \MEAH and \MEACOO, and trace-ion quantities retained from literature-backed or diagnostic calculations.  This hierarchy is used to evaluate a fixed advanced-Born parameter set for \COtwo-solubility modeling; it is not a pressure-only fitted model.

\begin{table}[t]
    \centering
    \caption{Full ionic ePC-SAFT component parameters used for the promoted SSM+DS MEA evaluation.}
    \label{tab:full-ionic-ssm-ds-parameters}
    \scriptsize
    \setlength{\tabcolsep}{3pt}
    \begin{tabularx}{\linewidth}{lrrrrrrr>{\raggedright\arraybackslash}X}
        \toprule
        Species & $m$ & $\sigma$ & $\epsilon/k$ & $z$ & $\varepsilon_r$ & $d^{\mathrm{Born}}$ & $f_k$ & Basis \\
        \midrule
        CO2 & 2.079 & 2.7852 & 169.21 & 0 & 1.4122 & 0 & 1 & Inherited neutral molecular parameter. \\
        MEA & 3.0353 & 3.0435 & 277.17 & 0 & 32 & 0 & 1 & Inherited associating amine; no adjusted MEA-specific $f_k$. \\
        H2O & 1.2047 & \makecell[l]{$2.7927 + 10.11\,e^{-0.01775T}$\\$- 1.417\,e^{-0.01146T}$} & 353.95 & 0 & 78.09 & 0 & 1.5 & Literature water solvent; $f_k=1.5$. \\
        MEAH+ & 1 & 3.4851 & 232.69 & 1 & 8 & 3.5332 & 1 & Fitted to Tier A MEA speciation. \\
        MEACOO- & 1 & 3.5354 & 453.27 & -1 & 8 & 3.5411 & 1 & Fitted to Tier A MEA speciation; corrected MW. \\
        HCO3- & 1 & 2.9296 & 70 & -1 & 8 & 3 & 1 & Held/Uyan bicarbonate values; regularized trace fit retained 3.0. \\
        CO3\textasciicircum{}2- & 1 & 2.4422 & 249.26 & -2 & 8 & 3 & 1 & Held/Uyan carbonate values; regularized trace fit retained 3.0. \\
        H3O+ & 1 & 3.4654 & 500 & 1 & 8 & 1.218 & 1 & Held/Uyan proton carrier with Figiel SSM+DS Born radius. \\
        OH- & 1 & 2.0177 & 650 & -1 & 8 & 3.0811 & 1 & Held/Uyan hydroxide values; hydration-derived Born diameter. \\
        \bottomrule
    \end{tabularx}
\end{table}


The parameter-evidence matrix in \cref{tab:parameter-evidence-matrix,tab:trace-parameter-evidence} distinguishes the evidentiary status of the retained inputs from trace-species quantities that are not independently identified by the available NMR targets.  The available speciation data test \MEAH and \MEACOO most directly; bicarbonate, carbonate, hydronium, and hydroxide quantities require additional caution because they are constrained by fewer direct targets or by balance relationships.

\begin{table*}[t]
    \centering
    \scriptsize
    \setlength{\tabcolsep}{3pt}
    \renewcommand{\arraystretch}{1.08}
    \caption{Parameter-evidence matrix for the primary true-species \MEA ePC-SAFT basis.}
    \label{tab:parameter-evidence-matrix}
    \begin{tabularx}{\textwidth}{>{\raggedright\arraybackslash}p{1.5cm}>{\raggedright\arraybackslash}p{3.1cm}>{\raggedright\arraybackslash}p{3.2cm}X}
        \toprule
        Species & Parameter scope & Evidence status & Data or basis \\
        \midrule
        \COtwo & \makecell[l]{$m$, $\sigma$, $\epsilon/k$,\\$\varepsilon_r$, MW} & Molecular PC-SAFT value & Non-electrolyte \COtwo parameterization used as the inherited molecular component. \\
        \MEA & \makecell[l]{$m$, $\sigma$, $\epsilon/k$,\\association terms, MW} & Molecular PC-SAFT value & Associating neutral amine parameterization retained as the \MEA molecular basis. \\
        \MEA & $f_k$ & Not adjusted in this work & Figiel et al. report solvent $f_k$ values near 1.5 for water and alcohol solvents, but the adopted \MEA calculation keeps $f_k=1$ until direct \MEA dielectric and activity evidence supports changing it. \\
        \HtwoO & \makecell[l]{$m$, $\sigma(T)$, $\epsilon/k$,\\association terms,\\$\varepsilon_r$, $f_k$, MW} & SSM+DS solvent parameter & Water PC-SAFT parameters and SSM+DS $f_k=1.5$ from modified-Born ePC-SAFT literature. \\
        \MEAH & \makecell[l]{$\sigma$, $\epsilon/k$,\\$d^{\mathrm{Born}}$} & Fitted to \MEA speciation data & Direct \MEAH NMR/speciation targets in the selected amine-ion data subset. \\
        \MEAH & \makecell[l]{$\varepsilon_r$, $f_k$,\\$z$, MW} & Ionic conventions & All ions use relative dielectric constant 8 and $f_k=1$ in SSM+DS; charge and molecular weight come from the species definition. \\
        \MEACOO & \makecell[l]{$\sigma$, $\epsilon/k$,\\$d^{\mathrm{Born}}$,\\$k_{ij}(\MEAH,\MEACOO)$} & Fitted to \MEA speciation data & Direct carbamate NMR/speciation targets in the selected amine-ion data subset; the \MEAH/\MEACOO interaction was fitted with the same subset. \\
        \MEACOO & \makecell[l]{$\varepsilon_r$, $f_k$,\\$z$, MW} & Ionic conventions plus corrected molecular weight & All ions use relative dielectric constant 8 and $f_k=1$ in SSM+DS; MW is corrected to \MEA+\COtwo-H = 0.10408 kg/mol. \\
        \bottomrule
    \end{tabularx}
\end{table*}

\begin{table*}[t]
    \centering
    \scriptsize
    \setlength{\tabcolsep}{3pt}
    \renewcommand{\arraystretch}{1.08}
    \caption{Trace-species parameter evidence retained for the true-species \MEA ePC-SAFT basis.}
    \label{tab:trace-parameter-evidence}
    \begin{tabularx}{\textwidth}{>{\raggedright\arraybackslash}p{1.5cm}>{\raggedright\arraybackslash}p{3.1cm}>{\raggedright\arraybackslash}p{3.2cm}X}
        \toprule
        Species & Parameter scope & Evidence status & Data or basis \\
        \midrule
        \HCOthree & \makecell[l]{$\sigma$, $\epsilon/k$,\\$d^{\mathrm{Born}}$,\\$\varepsilon_r$, $f_k$,\\$z$, MW} & Literature-backed trace-species set plus Born sensitivity & \HCOthree uses Held/Uyan $\sigma=2.9296$, $\epsilon/k=70.00$, and water-ion $k_{ij}=0.0$; the regularized calculation retains $d^{\mathrm{Born}}=3.0$, while an unanchored sensitivity check gives a lower-residual bicarbonate-sensitive value near 6.80. \\
        \COthree & \makecell[l]{$\sigma$, $\epsilon/k$,\\$d^{\mathrm{Born}}$,\\$\varepsilon_r$, $f_k$,\\$z$, MW} & Literature-backed trace-species set plus Born sensitivity & \COthree uses Held/Uyan $\sigma=2.4422$, $\epsilon/k=249.26$, and water-ion $k_{ij}=-0.25$; the regularized calculation retains $d^{\mathrm{Born}}=3.0$. \\
        \HthreeO & \makecell[l]{$\sigma$, $\epsilon/k$,\\$d^{\mathrm{Born}}$,\\$\varepsilon_r$, $f_k$,\\$z$, MW} & Literature-backed proton-carrier value & \HthreeO represents the proton carrier for acid--base equilibria; it uses Held/Uyan $\sigma=3.4654$, $\epsilon/k=500.00$, water-ion $k_{ij}=0.25$, and Figiel $d^{\mathrm{Born}}=1.218$. \\
        \Hydroxide & \makecell[l]{$\sigma$, $\epsilon/k$,\\$d^{\mathrm{Born}}$,\\$\varepsilon_r$, $f_k$,\\$z$, MW} & Literature-backed hydroxide value plus hydration estimate & \Hydroxide uses Held/Uyan $\sigma=2.0177$, $\epsilon/k=650.00$, and water-ion $k_{ij}=-0.25$; $d^{\mathrm{Born}}=3.0811$ is derived by Born hydration-energy inversion and is not an \MEA-system regression. \\
        \bottomrule
    \end{tabularx}
\end{table*}


\subsection{Pure-Component and Binary Parameters}
The molecular \COtwo, \MEA, and \HtwoO parameters are inherited from the cited PC-SAFT and ePC-SAFT sources listed in \cref{tab:full-ionic-ssm-ds-parameters}.  The temperature-dependent water segment diameter is
\begin{equation}
    \sigma_{\HtwoO}(T)=2.7927+10.11\exp(-0.01775T)-1.417\exp(-0.01146T),
    \label{eq:water-sigma-temperature}
\end{equation}
where \(T\) is in kelvin and \(\sigma_{\HtwoO}\) is in angstroms.  The final column of \cref{tab:full-ionic-ssm-ds-parameters} reports a source or status for every retained value; provisional inputs are explicitly distinguished from literature values.

The binary-interaction matrix is intentionally sparse.  \Cref{tab:binary-interaction-parameters} lists the nonzero interactions, the zero-valued candidate interactions retained for the fixed activity-based evaluation, and the transferred water--ion terms.  The remaining \MEA-specific ion--molecule and ion--ion terms are held at zero.  This sparse interaction structure tests the selected activity model before introducing a broader optimized binary-interaction matrix.

\begin{table}[t]
    \centering
    \caption{Nonzero binary dispersion interaction parameters used in the activity-based ePC-SAFT calculation.  Unlisted binary pairs use \(k_{ij}=0\).}
    \label{tab:binary-interaction-parameters}
    \small\rmfamily
    \setlength{\tabcolsep}{5pt}
    \renewcommand{\arraystretch}{1.08}
    \begin{tabularx}{\linewidth}{>{\raggedright\arraybackslash}p{0.26\linewidth}c>{\raggedright\arraybackslash}X}
        \toprule
        Pair & \(k_{ij}\) & Source \\
        \midrule
        \MEA--\HtwoO & -0.0520 & \cite{Baygi2015,Nasrifar2010} \\
        \MEAH--\MEACOO & -0.002018 & This work \\
        \HtwoO--\HthreeO & 0.2500 & \cite{Held2014,Uyan2015} \\
        \HtwoO--\Hydroxide & -0.2500 & \cite{Held2014,Uyan2015} \\
        \HtwoO--\COthree & -0.2500 & \cite{Held2014,Uyan2015} \\
        \bottomrule
    \end{tabularx}
\end{table}


\subsection{Fixed parameter evidence boundary}
The coupled native regression was not completed.  The retained \MEAH and \MEACOO values are provisional fixed inputs inherited from an exploratory historical calculation, not a reproducible parameter fit executed by the current workflow.  The NMR/speciation studies\cite{Matin2012,Jakobsen2005,Bottinger2008} are therefore used only to evaluate the resulting fixed-parameter speciation residuals.  Historical local-fit bounds, improvement metrics, and parity plots are excluded because they do not establish a current parameter-estimation result.  A future parameter claim requires an executed package-native, multi-observable regression with complete row provenance and held-out evidence defined before fitting.

\subsection{Evaluation residuals and future coupled objective}
The pressure residual is expressed on a logarithmic basis,
\begin{equation}
    \epsilon_{P,j}
    =
    \log_{10}\left(\frac{P_{\COtwo,j}^{\mathrm{model}}}{P_{\COtwo,j}^{\mathrm{exp}}}\right),
    \label{eq:pressure-log-residual}
\end{equation}
because measured \COtwo partial pressures span several orders of magnitude across the solubility data.  Speciation residuals use the same \(\log_{10}(\mathrm{model}/\mathrm{data})\) form for nonzero measured targets.  Measurements reported as zero and treated as upper-bound targets are evaluated separately with modeled mole-fraction upper bounds, and quantities inferred from balance constraints are used as context rather than direct logarithmic targets.

For a future jointly fitted model, the preregistered multi-observable target would combine pressure residuals, speciation residuals, reaction-equilibrium residuals, and regularization to the parameter seed:
\begin{equation}
    \begin{aligned}
    \Phi(\boldsymbol{\theta}) ={}&
    \sum_{j\in\mathcal{D}_P} w_{P,j}\epsilon_{P,j}^{2} \\
    &+
    \sum_{j\in\mathcal{D}_x}\sum_{i\in\mathcal{S}_x} w_{x,ij}
    \left[\log_{10}\left(\frac{x_{ij}^{\mathrm{model}}}{x_{ij}^{\mathrm{exp}}}\right)\right]^2 \\
    &+
    \sum_{j\in\mathcal{D}_x}\sum_r w_{K,rj}
    \left(\sum_i \nu_{ri}\ln a_{ij}-\ln K_r(T_j)\right)^2 \\
    &+
    \lambda\lVert \boldsymbol{\theta}-\boldsymbol{\theta}_0\rVert^2 .
    \end{aligned}
    \label{eq:regression-objective}
\end{equation}
Equation~\eqref{eq:regression-objective} is not minimized in the present work and supplies no fitted-parameter claim.  The \(K_r(T)\) values are the fixed reaction constants listed in \cref{tab:reaction-equilibrium-constants}.  A solubility-pressure-only fit can hide a poor liquid composition, while a speciation-only fit cannot guarantee the correct vapor-phase pressure.  A future thermodynamic regression intended for reactive \COtwo solubility and liquid speciation must therefore satisfy both observable families.

\subsection{Computational implementation and reproducibility}
The ideal baseline solves eight reaction, balance, and charge residuals in an eight-variable softmax representation using SciPy nonlinear least squares.  The first loading uses a chemistry-informed positive composition; subsequent points at the same temperature use the previous accepted composition as a continuation start.  The solve permits at most 1000 residual evaluations, uses \(10^{-11}\) function, step, and gradient tolerances, and is accepted only when the largest unscaled residual is at most \(10^{-8}\).

The activity-based calculation uses the same ordered nine-species basis with the pinned \texttt{epcsaft} 1.5.2 source at Git commit \texttt{9f51afd0f9c11a6497ddca05c8b2dd0ea0ffa785}.  Native reactive-speciation calls use at most 80 iterations, a \(10^{-7}\) solver tolerance, damping 0.5, and a minimum mole fraction of \(10^{-14}\).  Acceptance additionally requires mass-balance residuals no larger than \(10^{-7}\), charge residual no larger than \(10^{-6}\), reaction residuals no larger than \(10^{-7}\), finite positive normalized compositions, and zero state-evaluation failures.  Reactive bubble calculations use at most 120 iterations, a \(10^{-6}\) fugacity tolerance, and pressure bounds from \(10^{-3}\) to \(10^{8}\) Pa.  Best-effort solver returns are retained as diagnostics but are excluded from metrics and plots unless every acceptance condition passes.

The present work does not execute or report a coupled solubility/speciation regression.  The reported parameter set is evaluated as a fixed activity-based parameter set with explicit pressure and speciation residuals; no optimizer settings or runtime are therefore claimed for parameter estimation.  \Cref{tab:residual-summary} summarizes the residual metrics that define the manuscript claims.

\begin{center}
    \begin{minipage}{\linewidth}
    \centering
    \scriptsize
    \captionsetup{hypcap=false}
    \captionof{table}{Controlled pressure comparison and separate full-set/speciation context for the fixed-parameter \COtwo-solubility assessment.  Pressure residuals are \(\log_{10}(P_{\mathrm{model}}/P_{\mathrm{data}})\).}
    \label{tab:residual-summary}
    \begin{adjustbox}{max width=\linewidth}
    \begin{tabular}{llcl}
        \toprule
        Scope & Calculation & Rows & Residual summary \\
        \midrule
        Same-record & Ideal baseline & 31 & Median absolute 0.160; signed median \(-0.024\); RMSE 0.297 \\
        Same-record & Activity-based & 31 & Median absolute 0.495; signed median 0.226; RMSE 0.565 \\
        Full-set context & Activity-based & 161 & Median absolute 0.493; signed median \(-0.288\); RMSE 0.645 \\
        Speciation evidence & Ideal baseline, positive measurements & 74 states & \MEA 0.055; \MEAH 0.044; \MEACOO 0.029; \HCOthree 0.227; \ce{MEA + MEAH+} 0.0087 \\
        Speciation evidence & Activity-based, positive measurements & 74 states & \MEA 0.074; \MEAH 0.033; \MEACOO 0.040; \HCOthree 0.461; \COthree 0.599; \ce{MEA + MEAH+} 0.066 \\
        Speciation evidence & Source-reported zero \HCOthree values & 15 & Maximum modeled mole fraction \(=1.28\times10^{-3}\); balance-derived quantities are excluded from independent residuals \\
        \bottomrule
    \end{tabular}
    \end{adjustbox}
    \end{minipage}
\end{center}

